\section{问题重述}

\subsection{问题背景}

以大语言模型（LLM）为代表的人工智能技术正沿着“扩大规模”的路径快速演进。经典标度律\cite{kaplan2020scaling,hoffmann2022compute}表明，验证集交叉熵损失随参数量 $N$ 与训练 Token 数 $D$ 按幂律下降，训练算力近似为 $C\approx 6ND$。然而，单纯扩大规模正面临两重约束：一是算力增长趋缓，二是可用高质量语料日益稀缺，即所谓“数据墙”\cite{villalobos2024rundata,muennighoff2023dataconstrained}。与此同时，数据质量与领域配比对模型能力的作用日益凸显：仅 80 亿参数的 OpenScholar-8B 在引文准确率等任务上优于 GPT-4o\cite{asai2026openscholar}，以“教科书级”数据训练的小模型同样表现出超越规模的能力\cite{gunasekar2023textbooks}。

提升质量、扩大规模、延长上下文都要消耗算力，而调整领域配比虽不增加总开销，却改变算力在各领域间的分配效果。因此，在有限算力下如何在参数规模 $N$、数据规模 $D$、数据质量 $Q$、领域配比 $\boldsymbol p$ 与上下文长度 $L_{\mathrm{ctx}}$ 之间取舍，是一个典型的多资源约束优化问题。本题要求利用真实公开数据与已标注的半合成补充数据，刻画上述因素与 Loss 的定量关系，求解算力约束下的最优资源配置，并据此预测开源大模型的能力前沿。

\subsection{需要解决的问题}

四个问题构成“数据属性—损失规律—资源决策—能力前沿”的递进主线，前一问的输出是后一问的输入。

\textbf{问题一：数据质量评价、质量冲突消解与领域配比建模。}（1）对附件 A1--A3 的 22 项质量指标统一方向并预处理，建立综合评价模型，给出样本级与领域级质量评分 $Q$，说明聚合方式，并将全量结果与抽样集 A1 对照；（2）定义不同指标对同一文本评价显著不一致的“质量冲突”，分析成因并建立消解规则，在抽样集上分析、在扩展集上检验；（3）利用 A4--A15 建立 17 域配比 $\boldsymbol p$ 与交叉熵损失的定量关系，处理单纯形约束，在检验集 A6--A11 上验证，并用外推表 A12--A15 讨论外推稳健性。

\textbf{问题二：跨维度数据融合与广义标度律。}在经典标度律 $L=E+AN^{-\alpha}+BD^{-\beta}$ 基础上，综合附件 A、B 建立同时包含 $N,D,Q,\boldsymbol p$ 的广义标度律，提出统一两类实验的可检验假设；完成参数估计、检验与验证（B1 为主拟合，B2/B3、B4/B5 做族外与文献验证，B6--B8 为质量数据，B9/B10 讨论百亿参数以上外推）；分析各因素的边际效用与弹性、领域间的替代与互补关系，并推导质量提升与规模扩大可相互替代的条件。

\textbf{问题三：算力约束下的多维资源联合优化与结构性转移。}总算力由基础训练 $6ND$、质量提升 $D[g(Q)-g(Q_0)]_+$ 与长文本注意力 $\eta NDL_{\mathrm{ctx}}$ 三部分构成（$\eta=2\times10^{-4}$，$L_{\mathrm{ctx}}$ 依据 C7 外生给定）。在预算 $C$ 下求解最优资源配置，比较三类质量成本函数的影响；给出注意力开销与训练开销相当的临界长度 $L_{\mathrm{ctx}}^{\mathrm{crit}}=6/\eta$ 并做敏感性分析；考察 $C=10^{19},10^{22},10^{24}$ FLOPs 等不同量级预算下最优策略是否发生结构性转移，给出其数学定义与识别方法。

\textbf{问题四：技术演进分析与前沿预测。}结合附件 C 建立动力学或因果推断模型，分离并量化规模扩张与非规模技术进步对能力提升的贡献占比；利用 Loss--Benchmark 桥接数据把前三问的 Loss 结论映射到下游能力并讨论映射误差；说明综合能力度量、开源筛选口径、模型类型与时间轴口径，在算力增长放缓情景下预测未来 12 或 24 个月开源大模型能力前沿并给出不确定性分析。
